\section{Results}
\label{sec:results}

\subsection{Prediction Accuracy}

Table~\ref{tab:regression} reports regression performance on the 18-game validation sample.
The TIGER model substantially outperforms the weight-only baseline on both metrics,
\textbf{exceeding both pre-registered thresholds}.

\begin{table}[ht]
\centering
\caption{Regression model performance ($N = 18$ games, BGG community ratings as ground truth).}
\label{tab:regression}
\begin{tabular}{lrrl}
\toprule
\textbf{Feature set} & \textbf{Pearson $r$} & \textbf{RMSE} & \textbf{vs.\ baseline} \\
\midrule
F1: Weight-only (baseline) & 0.187 & 0.636 & --- \\
F2: TIGER 5-dim only       & \textbf{0.715} & \textbf{0.453} & \textbf{+28.8\% RMSE reduction} \\
\bottomrule
\multicolumn{4}{l}{\footnotesize Pre-registered thresholds: $r \geq 0.60$; $\geq 15\%$ RMSE reduction.} \\
\multicolumn{4}{l}{\footnotesize Both thresholds passed with margin.} \\
\end{tabular}
\end{table}

The weight-only model ($r = 0.187$) performs near-chance, confirming that BGG
complexity weight is a poor predictor of community satisfaction --- a finding that
motivates multi-dimensional profiling.
The TIGER model reaches $r = 0.715$ and reduces RMSE by 28.8\%,
demonstrating that the five TIGER dimensions capture player-relevant variation
beyond game complexity.

\textbf{Key failure case:} Catan ($\hat{y} = 8.00$, actual $= 7.10$) is over-predicted
because its TIGER profile (I=6.0, moderate on all dimensions) resembles higher-rated
games. This is consistent with Catan's mixed community reception: it invented modern
hobby gaming but is widely considered mechanically outdated.
Chess ($\hat{y} = 6.66$, actual $= 5.5$) is over-predicted for analogous reasons.
Both cases suggest a \emph{historical legacy} factor TIGER does not capture.

\textbf{Key success cases:} Brass: Birmingham ($\hat{y}=8.19$, actual=8.6) and
Spirit Island ($\hat{y}=8.03$, actual=8.5) --- the two highest-rated games ---
are correctly predicted in the top tier.

\subsection{Ablation Analysis}

Table~\ref{tab:ablation} shows how Pearson $r$ changes when each TIGER dimension
is dropped from the model.

\begin{table}[ht]
\centering
\caption{Drop-one ablation: $r$ when each dimension is excluded.}
\label{tab:ablation}
\begin{tabular}{lrr}
\toprule
\textbf{Dimension dropped} & \textbf{$r$ (4-dim)} & \textbf{$\Delta r$} \\
\midrule
T (Tension)     & 0.553 & $-$0.162 \hspace{2pt}\textbf{(largest loss)} \\
E (Elegance)    & 0.613 & $-$0.102 \\
I (Interaction) & 0.707 & $-$0.008 \\
R (Range)       & 0.704 & $-$0.011 \\
G (Architecture)& 0.715 & $-$0.001 \\
\midrule
Full (T,I,G,E,R) & \textbf{0.715} & --- \\
\bottomrule
\end{tabular}
\end{table}

\textbf{Tension (T) is the single most predictive dimension} ($\Delta r = 0.162$).
This is counterintuitive: higher catastrophe pressure predicts higher BGG ratings.
Interpretation: games with escalating threat mechanics (Spirit Island, Agricola, Brass)
tend to attract committed hobbyists who rate generously, while low-tension light games
(Crokinole, Skull) attract casual players who rate lower on the BGG scale.
This is a \emph{selection effect}: T predicts not gameplay quality but the type of
player community that forms around the game.

Elegance (E) contributes $\Delta r = 0.102$.
The negative coefficient (higher elegance predicts lower BGG ratings)
confirms that accessibility and community rating are negatively correlated at the top
of the BGG scale: the highest-rated games (Brass Birmingham, Spirit Island) are complex
and hard to teach; the most elegant games (Crokinole, Skull) rate modestly.

I, G, and R contribute minimally in this sample ($|\Delta r| \leq 0.011$),
suggesting their predictive signal is largely captured by T and E.
This may reflect sample composition: the 18-game sample under-represents the
mid-range games where I and G vary most independently.
